\section{Парадокс Рассела}

\begin{axiom}{3.8}{универсальная спецификация}
Для любого утверждения $P$ существует множество $A$ такое, что для всякого объекта $x$
\[
  x \in A \iff P(x).
\]
\leansource{axiom_of_universal_specification}{: Prop := ∀ P : Object → Prop, ∃ A : Set, ∀ x : Object, x ∈ A ↔ P x}
\end{axiom}

Эта аксиома кажется естественным способом ввести в теорию произвольные множества, задаваемые
свойством. Однако она противоречива.

\begin{theorem*}{парадокс Рассела}
Аксиома универсальной спецификации ложна.
\leansource{Russells_paradox}{: ¬ axiom\_of\_universal\_specification}
\end{theorem*}

\begin{proof}
Рассмотрим свойство
\[
  P(x) :\iff \exists X.\; (x = X \land x \notin X),
\]
то есть «$x$ является множеством, не содержащим самого себя». По аксиоме универсальной
спецификации, применённой к $P$, существует множество $\Omega$ такое, что для любого $x$
\[
  x \in \Omega \iff P(x).
\]
Применим это к $x = \Omega$ и разберём два случая.

\caselabel{$\Omega \in \Omega$} Тогда $P(\Omega)$ верно, то есть найдётся множество
$X$ с $\Omega = X$ и $\Omega \notin X$. Подставляя $X = \Omega$, получаем $\Omega \notin \Omega$ —
противоречие с предположением случая.

\caselabel{$\Omega \notin \Omega$} Тогда само $\Omega$ свидетельствует об истинности
$P(\Omega)$ (берём $X := \Omega$), а значит $\Omega \in \Omega$ по построению $\Omega$ —
противоречие.

В обоих случаях приходим к противоречию, значит исходное предположение (аксиома универсальной
спецификации) ложно.
\end{proof}

\begin{remark}
Парадокс показывает, что «множество всех множеств, не содержащих себя» не может существовать —
классическая формулировка парадокса Рассела. Именно поэтому в аксиоматической теории множеств
вместо неограниченной спецификации используется \emph{ограниченная} спецификация: $P$
применяется только к элементам уже существующего множества.
\end{remark}

Чтобы дополнительно исключить патологии вроде $\Omega \in \Omega$, вводится ещё одна аксиома.

\begin{axiom}{3.9}{регулярность}
Всякое непустое множество $A$ содержит элемент $x$, который либо не является множеством, либо,
будучи множеством $S$, не пересекается с $A$:
\[
  \exists x \in A.\; \forall S.\; (x = S \implies S \cap A = \varnothing).
\]
\leansource{SetTheory.Set.axiom_of_regularity}{(h : A ≠ ∅) : ∃ x : A, ∀ S : Set, x.val = S → Disjoint S A}
\end{axiom}

\begin{remark}
Аксиома регулярности запрещает бесконечный спуск по цепочке принадлежности
$\cdots \in x_2 \in x_1 \in x_0$ и, в частности, запрещает множеству содержать само себя
($A \in A$ невозможно).
\end{remark}
